\subsection{4.2  Phase 2: Cause Attribution}

Once triggered, the APIP requires a structured cause attribution process before any intervention is applied. The goal is to identify which failure mode category (Section 3) best explains the degradation. Tier mismatch from misdiagnosis is the primary source of avoidable regression: prompt engineering applied to a retrieval freshness problem, for instance, consumes engineering cycles with no improvement.

This attribution problem is distinct from, and complementary to, recent work on automated failure attribution in agentic systems. That literature localizes the responsible agent or step within a single failed trajectory. Causal-inference frameworks combine Shapley-based agent-level blame assignment with causal discovery over execution logs to identify critical failure steps, motivated by the finding that state-of-the-art methods achieve under 15\% accuracy at locating root-cause steps on the Who\&When benchmark [36, 37]. Lightweight causal-graph tracing supports post-hoc root-cause localization in deployed multi-agent workflows [38]. APIP Phase 2 operates one level of aggregation above: it classifies a window-level degradation into a remediation-oriented failure-mode category to select an intervention tier. The two levels compose rather than compete. Trajectory-level attribution outputs are natural evidence inputs to Phase 2, since recurring step- or agent-level attributions across sampled failures constrain the failure locus and the plausible failure-mode classification. Recent root-cause taxonomies make the point explicitly: localizing the failing step falls short of diagnosing the underlying cause, whether poor prompt design, improper task decomposition, or logical deadlock [39]. That diagnostic gap is what the APIP's mandatory taxonomy classification is designed to close at the deployment level.

The cause attribution process should include: (1) disaggregated performance analysis by input category, user segment, and topic cluster to isolate failure locus; (2) analysis of a sampled set of failure cases, applying automated trajectory-level attribution methods where tooling permits [36, 38], with a qualified human reviewer adjudicating the resulting evidence; (3) retrieval audit if behavioral drift is suspected; (4) tool invocation log analysis if tool misuse is suspected; and (5) a structured classification of the failure into the four-category taxonomy. The phase closes with a written cause attribution report recording the evidence examined, any automated attribution outputs, and the adjudicated classification; this report becomes part of the APIP record.

